\section{Transport along the leaf, and the second branch}\label{sec:leaf}

Section \ref{sec:point} established, structurally, that exactly one Painlev\'e
V leaf passes through System $\mathsf C$, and \S\ref{sec:monodromy} produced
the point of the character variety that the system represents. Neither of those
is yet a demonstration that the system \emph{moves along} that leaf. This
section performs the test: integrate the flow in $t$, rebuild the differential
system at the new time, recompute the trace coordinates, and ask whether they
have moved.

They should not. The trace coordinates are the isomonodromic invariants; a
deformation is isomonodromic precisely when they are constant. So the
observable is the drift, and the test is passed by its absence --- which is why
\S\ref{ssec:sensitivity} had to establish first that the drift is not absent
for trivial reasons.

\subsection{The transport test}\label{ssec:transport}

\begin{remark}[Constancy along the flow; \texttt{S22-15}, \textbf{VERIFIED}]
System $\mathsf C$'s point transports along Painlev\'e V with the trace
coordinates $x_1,x_2,x_3$ constant to $90.81$ digits at $\varepsilon=10^{-3}$
and $91.83$ digits at $\varepsilon=10^{-4}$, against a pre-registered threshold
of $30$ digits. The gap left open in the location analysis is thereby closed
\emph{numerically}. This is a numerical agreement at one point along one
direction in $t$; \emph{it is not a proof of leaf membership}.
\end{remark}

The last sentence is part of the claim, not a hedge added in drafting, and the
label is \textbf{VERIFIED} for exactly that reason. Ninety digits is a great
deal of agreement and it is still agreement at a point. What would be required
for more is set out in \S\ref{ssec:membership-verdict}; what is available is
this.

A constancy result is only as good as its ability to fail, so the same machine
was run with the flow deliberately broken.

\begin{remark}[The test discriminates; \texttt{S22-16}, \textbf{VERIFIED}]
The same transport run with a deliberately wrong vector field --- the constant
$2\theta_\infty-2$ dropped from $dp/dt$ --- moves the trace coordinates by
$3.83\times10^{-2}$.
\end{remark}

The mutation is small and structural: one constant term removed from one
component of the field \eqref{eq:flow}, leaving a vector field that is still a
plausible flow and still integrates without complaint. It moves the
coordinates by nearly ninety orders of magnitude more than the true field does.
The constancy in \texttt{S22-15} is therefore a property of the specific flow
of \S\ref{ssec:flow} and not of the integration scheme, the rebuilding
procedure, or the coordinates.

\subsection{The residual is not round-off}\label{ssec:notroundoff}

Ninety digits of agreement out of a working precision of two hundred invites
the question of what limits it. If the answer were arithmetic precision, the
result would be an artefact of the arithmetic; the claim is that it is not.

\begin{remark}[Precision is not the ceiling; \texttt{S22-17},
\textbf{VERIFIED}]
Raising the precision from $200$ to $320$ digits leaves the agreement and its
ratio unchanged to every printed figure --- $90.8118$, $91.8301$, ratio
$10.429999$ --- while every precision-limited control improves by about $120$
digits: $\|MMM-I\|$ from $1.20\times10^{-105}$ to $5.57\times10^{-225}$; the
cubic from $3.96\times10^{-112}$ to $7.13\times10^{-232}$; $|x_3(t_1)-x_3(t_0)|$
from $2.35\times10^{-113}$ to $3.30\times10^{-233}$. The $\alpha$ residual stays
at $1.834\times10^{-91}$. Over $\varepsilon=10^{-3},10^{-4},10^{-5}$ the ratios
are $10.430$ and $10.034$. The residual is therefore a systematic $O(\varepsilon)$
error and not round-off; it is floored by the optimal-truncation level
$e^{-|t|R}=10^{-88.7}$ set by the pre-registered $R$, which no precision can
reduce.
\end{remark}

This is the cleanest structural statement in the section, and it is worth
saying why. Two families of quantities were measured at two precisions. One
family moved by $120$ digits and one did not move at all. A quantity that is
insensitive to a $120$-digit change in precision is not limited by precision.
The ratios near $10$ identify what it \emph{is} limited by --- first-order
dependence on the step $\varepsilon$ --- and the optimal-truncation bound of
\texttt{S22-09} identifies the floor beneath that. The three measurements
agree, and the agreement is what converts ``the residual is small'' into ``the
residual is understood''.

\subsection{The second branch}\label{ssec:branch2}

The location analysis produced a \emph{pair} of solutions, conjugate in
$\sqrt3$. Everything above concerns one of them. The second was not inferred
from the first.

\begin{remark}[Branch 2, run independently; \texttt{S22-18},
\textbf{VERIFIED}]
Branch $2$ of the sign pair, $t=+2\sqrt{33}\,i/9$, passes the convention
selection, the sensitivity check and the transport test independently: the
selection test picks exactly one convention out of $16$, this time
right${}|{}$clockwise${}|{}\gamma_\infty$ counterclockwise${}|{}01$; the
sensitivity ratios are $99.9952$ and $99.9998$; the transport agreement is
$90.82$ and $91.84$ digits against the same $30$-digit threshold. Branch $2$
was \emph{run}, not inferred.
\end{remark}

The emphasis is deliberate and \S\ref{ssec:galois} explains why it was
necessary rather than merely careful.

\subsection{The Galois relation between the branches}\label{ssec:galois}

\begin{remark}[The branches are Galois conjugate; \texttt{S22-19},
\textbf{PROVEN}]
Branches $1$ and $2$ are related by the automorphism $\sigma$ of
$\Q(\sqrt3,\sqrt{11})$ with $\sigma(\sqrt3)=-\sqrt3$ and
$\sigma(\sqrt{11})=+\sqrt{11}$, which maps $t$, $\theta_\infty$, $q$ and $p$
term by term, each difference simplifying to exactly zero, with an independent
numeric arm. They are \emph{not} complex conjugates.
\end{remark}

\begin{remark}[Why $\sigma$ does not transfer the monodromy; \texttt{S22-19},
\textbf{PROVEN}]
$\sigma$ is an automorphism of a number field. It does not act on analytic
continuation. It therefore transfers the location argument of
\S\ref{sec:point}, which is algebraic, and carries \emph{no} monodromy
information --- which is precisely why branch $2$ had to be run.
\end{remark}

This is the load-bearing distinction of the section. It is tempting to regard
two Galois-conjugate parameter values as interchangeable and to read off the
second branch's monodromy from the first. That inference is invalid: the
monodromy is computed by continuation along paths, and a field automorphism has
nothing to say about paths. The cost of the distinction is one complete
independent run; the benefit is that the statement below is a comparison of two
computed objects rather than a restatement of one.

\begin{remark}[The same point, up to relabelling; \texttt{S22-20},
\textbf{VERIFIED}]
The two branches yield the same point of the character variety up to
relabelling: $x_1$ and $x_2$ are interchanged, while $x_3$ and $\alpha$ are
identical, and the selected conventions are exact mirrors of one another. The
mirrored convention is a \emph{result} of branch $2$'s own selection test, not
an assumption carried over from branch $1$.
\end{remark}

\subsection{What the upgrade says, and what it does
not}\label{ssec:upgrade-scope}

The exhaustive elimination of \S\ref{ssec:exhaustive} changes the status of the
branch relation, and the change has to be stated with its scope attached.

\begin{remark}[The parameter-level upgrade; \texttt{S25-02},
\textbf{PROVEN}]
The two branches are Galois conjugate \emph{at the level of the parameters},
and this is exhibited rather than observed: the saturated matching ideal is
zero-dimensional of degree exactly $2$ over $K=\Q(i\sqrt{11})$, and its lex
eliminant in $\theta_\infty$ is proportional to
$27\theta_\infty^2+108\theta_\infty+107$, whose roots are $-2\pm\sqrt3/9$. The
map $\sqrt3\mapsto-\sqrt3$ is that eliminant's own Galois action.
\end{remark}

\begin{remark}[The scope of the upgrade]
\texttt{S25-02} supplies the structural content that \texttt{S22-20} lacked,
but it is a statement about the \emph{parameters}. It is not a statement about
the Stokes data, which is transcendental in the entries of the system. The
branch relation is therefore no longer purely numerical; it is \emph{not}
thereby proved at the level of the character variety. The two are different
statements and this paper does not elide the difference: \texttt{S22-20}
remains \textbf{VERIFIED}, on two independent runs.
\end{remark}

\subsection{An involution, and a structural route that does not
close}\label{ssec:involution}

If the branches give the same point up to an interchange of $x_1$ and $x_2$,
the natural question is whether some symmetry of the differential system
realises that interchange. There is an obvious candidate.

\begin{remark}[The involution; \texttt{S22-23}, \textbf{PROVEN}]
For the System $\mathsf C$ family at both branches the denominator
$D(z)=z^2-z$ is exactly invariant under $\varphi(z)=1-z$. Hence $\varphi$ fixes
the singular set $\{0,1,\infty\}$ setwise and exchanges the two finite
punctures. Since $\theta_0=\theta_1=0$, the local data at $0$ and at $1$ are
interchangeable, so $\varphi$ is the only candidate involution that could
produce the observed $x_1\leftrightarrow x_2$ swap with $x_3$ and $\alpha$
fixed.
\end{remark}

The candidate is natural, and it fails.

\begin{remark}[No constant-gauge equivalence; \texttt{S22-24},
\textbf{PROVEN}]
There is no analytic equivalence between the two branches of the form
$A_2(z)=-G\,A_1(\varphi(z))\,G^{-1}$ with $G$ a \emph{constant} gauge, for any
of the four candidate maps $\varphi$: the identity, $z\mapsto1-z$, coefficient
conjugation, and coefficient conjugation composed with $z\mapsto1-z$. Since $A$
is traceless, its determinant is the only gauge invariant, and the necessary
scalar condition on $-\det A$ fails for all four, with maximal coefficient
$3.83$, $0.416$, $0.416$ and $3.83$ respectively.
\end{remark}

\begin{remark}[The negative is informative]
A failed symbolic identity can mean ``no such relation'' or ``the simplifier
did not get there''. The positive control separates them: $\sigma$ applied to
branch $1$'s invariant, minus branch $2$'s, returns exactly zero through the
same radicals and the same simplifier. The machinery is therefore capable of
returning zero on this class of expression when zero is the answer, and the
four nonzero results are findings rather than failures to simplify.
\end{remark}

The consequence is stated plainly. The proposed structural route to
\texttt{S22-20} does not close, and \texttt{S22-20} stands as
\textbf{VERIFIED}-numerical on two independent runs, not as structural. What
has been refuted is the \emph{simplest form} of the argument: a $z$-dependent
gauge is untested and remains open, and is listed as such in
\S\ref{sec:open}.

This is the second place in the paper where a structural argument for a
numerically solid fact has been sought and not found --- the first being the
twist question of \S\ref{ssec:tracefree}, where one was found. Recording the
failure at the same weight as the success is the point: a reader who sees only
the successes cannot tell which of the paper's identifications are supported by
structure and which are supported by digits.
